\documentclass[11pt]{article}

\usepackage[T1]{fontenc}
\usepackage{lmodern}
\usepackage{microtype}
\usepackage[margin=1in]{geometry}
\usepackage{amsmath,amssymb,amsthm,mathtools}
\usepackage{booktabs,tabularx,array}
\usepackage{enumitem}

\newtheorem{proposition}{Proposition}
\newtheorem{lemma}[proposition]{Lemma}
\newtheorem{remark}[proposition]{Remark}

\newcommand{\R}{\mathsf R}
\newcommand{\D}{\mathsf D}
\newcommand{\Ssafe}{\mathsf S}
\newcommand{\Tr}{\operatorname{Tr}}
\newcommand{\Span}{\operatorname{span}}
\newcommand{\supp}{\operatorname{supp}}
\newcommand{\ip}[2]{\left\langle #1,#2\right\rangle}
\newcommand{\norm}[1]{\left\lVert #1\right\rVert}
\newcommand{\coeff}[2]{[z^{#1}]\,#2}

\title{An Excursion Interpretation of Proposition~5.4}
\author{}
\date{}

\begin{document}
\maketitle

\begin{abstract}
This note gives a word-by-word interpretation of the decomposition used in
Proposition~5.4.  Relative to the constant direction, the exceptional
eigendirection, and the remaining spectral subspace, every trace term is a
closed word in three letters.  A surviving word that visits the constant
direction decomposes uniquely into excursions from that direction.  The
one-excursion words give the terms $m k_m(\Lambda)\gamma^2$ and
$m\ip{g_s}{k_m(A_s)g_s}$.  The aggregate of all words with at least two
excursions is the coefficient of $L_-(z)-Y(z)$ and is nonnegative.  The two
words that never visit the constant direction give $-\Lambda^m$ and
$-\Tr(A_s^m)$.  A Hilbert--Schmidt estimate controls the latter and produces
the remaining $2M^{m-2}$ terms.  Combining these contributions gives exactly
Proposition~5.4.
\end{abstract}

\section{Setting and notation}

Let $W\colon[0,1]^2\to[0,1]$ be a symmetric step graphon, put
\[
  U=1-W,\qquad p=\int_{[0,1]^2}W,\qquad q=1-p,
\]
and let $T_U$ and $T_W$ be the corresponding integral operators.  We work in
the finite-dimensional space generated by the cells of the step graphon.  The
constant function $\mathbf 1$ has norm one.  Relative to
\[
  \mathbb R\mathbf 1\oplus E_0,
  \qquad E_0=\mathbf 1^\perp,
\]
write
\[
  T_U=
  \begin{pmatrix}
    q & g^*\\
    g & A
  \end{pmatrix},
  \qquad
  T_W=T_{\mathbf 1}-T_U=
  \begin{pmatrix}
    p & -g^*\\
    -g & -A
  \end{pmatrix}.
  \tag{1.1}
\]
Here
\[
  g=T_U\mathbf 1-q\mathbf 1\in E_0,
  \qquad A=P_{E_0}T_UP_{E_0}.
\]

Throughout the note, $m\geq3$ is odd.  In the regime relevant to
Proposition~5.4, $1/2<p<2/3$, hence $1/3<q<1/2$.  Suppose that $A$ has a
unique eigenvalue $\Lambda>q$, with a unit eigenfunction $\phi$.  Decompose
\[
  E_0=\mathbb R\phi\oplus H_s,
  \qquad
  g=\gamma\phi+g_s,
  \qquad
  \gamma=\ip{g}{\phi},
  \qquad g_s\perp\phi,
  \tag{1.2}
\]
and write $A_s=A|_{H_s}$.  Since $A$ is self-adjoint, $H_s$ is invariant
under $A$.  Relative to
\[
  \mathbb R\mathbf 1\oplus\mathbb R\phi\oplus H_s,
\]
the operator $T_W$ has the block form
\[
  T_W=
  \begin{pmatrix}
    p & -\gamma & -g_s^*\\
    -\gamma & -\Lambda & 0\\
    -g_s & 0 & -A_s
  \end{pmatrix}.
  \tag{1.3}
\]

Let $P_0,P_\Lambda,P_s$ denote the orthogonal projections onto
$\mathbb R\mathbf 1$, $\mathbb R\phi$, and $H_s$, respectively.  We label
these three states by
\[
  \R=P_0,
  \qquad
  \D=P_\Lambda,
  \qquad
  \Ssafe=P_s.
\]
The letters may be read as ``return state'', ``exceptional direction'', and
``remaining subspace''.  These labels are only mnemonic; all definitions
below are in terms of the projections.

\section{Closed projection words}

Insert $I=P_0+P_\Lambda+P_s$ between consecutive copies of $T_W$.  This gives
the exact expansion
\[
  \Tr(T_W^m)
  =\sum_{\omega\in\{\R,\D,\Ssafe\}^m}w(\omega),
  \tag{2.1}
\]
where, if $P_{\R}=P_0$, $P_{\D}=P_\Lambda$, and $P_{\Ssafe}=P_s$, then
\[
  w(\omega)
  :=\Tr\!\left(
    P_{\omega_1}T_WP_{\omega_2}T_W\cdots
    P_{\omega_m}T_W
  \right).
  \tag{2.2}
\]
The word is cyclic: $\omega_m$ and $\omega_1$ are adjacent because the final
copy of $T_W$ returns to the first projection inside the trace.

\subsection{Precise definition of an excursion}

For a word $\omega$, a \emph{visit to the constant direction} is an index
$i$ for which $\omega_i=\R$.  Suppose that $\omega$ has at least one such
visit and is not the constant word $\R^m$.  An \emph{excursion} is a maximal
nonempty cyclic interval of indices on which the letters lie in
$\{\D,\Ssafe\}$, with an $\R$ immediately before and immediately after the
interval.  Thus an excursion is a maximal block between two successive
visits to $P_0$.

The zero blocks in (1.3) give
\[
  P_\Lambda T_WP_s=P_sT_WP_\Lambda=0.
  \tag{2.3}
\]
Consequently, every word containing adjacent letters $\D\Ssafe$ or
$\Ssafe\D$ has zero weight.  It follows that every excursion in a nonzero
word is monochromatic: it has the form $\D^b$ or $\Ssafe^b$ for some
$b\geq1$.

If a word has no visit to $P_0$, then (2.3), including cyclic adjacency,
shows that it is either $\D^m$, $\Ssafe^m$, or a zero-weight word.  This is
the accurate meaning of the ``no return'' class; it is not an additional
restriction imposed on a walk.

\begin{center}
\renewcommand{\arraystretch}{1.25}
\begin{tabularx}{\textwidth}{@{}>{\raggedright\arraybackslash}p{0.19\textwidth}
  >{\raggedright\arraybackslash}p{0.22\textwidth}X@{}}
\toprule
Class of word & Representative & Exact contribution or conclusion\\
\midrule
Only the constant state
  & $\R^m$
  & $p^m$.\\
No visit to $P_0$; only the exceptional direction
  & $\D^m$
  & $(-\Lambda)^m=-\Lambda^m$.\\
No visit to $P_0$; only the remaining subspace
  & $\Ssafe^m$
  & $\Tr((-A_s)^m)=-\Tr(A_s^m)$.\\
No visit to $P_0$; both nonconstant states occur
  & for example $\D\D\Ssafe\Ssafe\D\cdots$
  & Zero, because some cyclically adjacent pair is $\D\Ssafe$ or
    $\Ssafe\D$.\\
Exactly one exceptional excursion of length $b$
  & $\R^{m-b}\D^b$, up to cyclic rotation
  & Each rotation has weight
    $p^{m-b-1}\gamma^2(-\Lambda)^{b-1}$.\\
Exactly one remaining-subspace excursion of length $b$
  & $\R^{m-b}\Ssafe^b$, up to cyclic rotation
  & Each rotation has weight
    $p^{m-b-1}\ip{g_s}{(-A_s)^{b-1}g_s}$.\\
At least two excursions
  & for example $\R\R\D\D\D\R\R\Ssafe\Ssafe\Ssafe\R$
  & Individual words need not be nonnegative.  Their aggregate is
    $m\coeff{m}{\sum_{r\ge 2} Y(z)^r/r}\geq0$.\\
\bottomrule
\end{tabularx}
\end{center}

For the displayed word of length $11$,
\[
  \omega=\R\R\D\D\D\R\R\Ssafe\Ssafe\Ssafe\R,
\]
there are three $\R$--$\R$ transitions, two exceptional internal
transitions, and two remaining-subspace internal transitions.  Its weight is
\[
  w(\omega)
  =p^3\gamma^2\Lambda^2\ip{g_s}{A_s^2g_s}.
  \tag{2.4}
\]
This example happens to be nonnegative.  In contrast, an exceptional
excursion of even length has an odd number of internal $-\Lambda$ factors and
therefore has negative weight.  Thus the assertion that some terms may be
discarded will be proved only after the appropriate words have been summed.

\section{The excursion generating function}

Define
\[
  Y(z):=
  \frac{z^2\ip{g}{(I+zA)^{-1}g}}{1-pz},
  \qquad
  L_-(z):=-\log(1-Y(z)).
  \tag{3.1}
\]
All series are formal power series.  Since $Y$ has order at least two,
\[
  L_-(z)=\sum_{r\geq1}\frac{Y(z)^r}{r}
  =Y(z)+\sum_{r\geq2}\frac{Y(z)^r}{r}
  \tag{3.2}
\]
is well defined coefficientwise.

The interpretation of a single factor $Y$ follows from
\[
  (I+zA)^{-1}=\sum_{j\geq0}z^j(-A)^j,
  \qquad
  \frac1{1-pz}=\sum_{k\geq0}p^kz^k.
\]
Hence
\[
  Y(z)
  =\sum_{j,k\geq0}
    z^{j+k+2}p^k\ip{g}{(-A)^jg}.
  \tag{3.3}
\]
For the term indexed by $(j,k)$:
\begin{itemize}[leftmargin=2em]
  \item the factor $z^2$ records the transition from $P_0$ to $E_0$ and the
        transition back to $P_0$;
  \item $\ip{g}{(-A)^jg}$ records the $j$ internal transitions of the
        nonconstant excursion;
  \item $p^k$ records the $k$ $P_0$--$P_0$ transitions outside that
        excursion.
\end{itemize}
Thus $Y$ is the generating function for one nonempty excursion together with
the constant steps that follow it.  A product $Y^r$ records an ordered list of
$r$ excursions.  In a cyclic word there are $r$ choices of which excursion
is declared first; the factor $1/r$ in (3.2) removes this choice.  Finally,
multiplication of the coefficient by $m$ roots the cyclic word at one of its
$m$ positions, as in the trace expansion (2.1).

The following determinant identity makes this interpretation exact, including
all multiplicities.

\begin{lemma}[Exact resummation of the excursion words]
For every odd $m\geq3$,
\[
  t(C_m,W)=\Tr(T_W^m)
  =p^m-\Tr(A^m)+m\coeff{m}{L_-(z)}.
  \tag{3.4}
\]
Here $p^m$ is the word $\R^m$, $-\Tr(A^m)$ is the aggregate of the words
with no visit to $P_0$, and $m\coeff{m}{L_-}$ is the aggregate of all words
that visit both $P_0$ and $E_0$.
\end{lemma}

\begin{proof}
From (1.1),
\[
  I-zT_W=
  \begin{pmatrix}
    1-pz & zg^*\\
    zg & I+zA
  \end{pmatrix}.
\]
Taking the Schur complement of $I+zA$ gives
\[
\begin{aligned}
  \det(I-zT_W)
  &=\det(I+zA)
    \left(1-pz-z^2\ip{g}{(I+zA)^{-1}g}\right)\\
  &=\det(I+zA)(1-pz)(1-Y(z)).
\end{aligned}
\]
Apply $-\log$ and use
\[
  -\log\det(I-zX)=\sum_{j\geq1}\frac{\Tr(X^j)}jz^j.
\]
The coefficient of $z^m$, multiplied by $m$, is
\[
  \Tr(T_W^m)=\Tr((-A)^m)+p^m+m\coeff{m}{L_-}.
\]
Since $m$ is odd, $\Tr((-A)^m)=-\Tr(A^m)$, proving (3.4).
\end{proof}

\section{Why the multi-excursion aggregate may be discarded}

Let
\[
  A=\sum_i\lambda_iP_i,
  \qquad
  \mu=\sum_i\norm{P_ig}^2\delta_{\lambda_i}.
  \tag{4.1}
\]
Thus
\[
  \ip{g}{F(A)g}=\int F(\lambda)\,d\mu(\lambda).
\]
The spectrum of $A$ lies in $[-1/2,1/2]$.  Since $p>1/2$, we have
$|\lambda|<p$ throughout the support of $\mu$.

\begin{lemma}[Nonnegative excursion coefficients]
Every coefficient of $Y(z)$ is nonnegative.  Consequently, every coefficient
of $L_-(z)-Y(z)$ is nonnegative.  More explicitly, for every $r\geq2$,
\[
  \coeff{r}{Y(z)}
  =\int
    \frac{p^{r-1}-(-\lambda)^{r-1}}{p+\lambda}
    \,d\mu(\lambda)
  \geq0.
  \tag{4.2}
\]
For odd $m$ this gives
\[
  \coeff{m}{L_-(z)}
  \geq\coeff{m}{Y(z)}
  =\int k_m(\lambda)\,d\mu(\lambda),
  \qquad
  k_m(\lambda):=\frac{p^{m-1}-\lambda^{m-1}}{p+\lambda}.
  \tag{4.3}
\]
\end{lemma}

\begin{proof}
Taking the coefficient of $z^r$ in (3.3) gives
\[
\begin{aligned}
  \coeff{r}{Y(z)}
  &=\int\sum_{j+k=r-2}p^k(-\lambda)^j\,d\mu(\lambda)\\
  &=\int
    \frac{p^{r-1}-(-\lambda)^{r-1}}{p+\lambda}
    \,d\mu(\lambda).
\end{aligned}
\]
The denominator is positive.  The numerator is positive because
$p>|\lambda|$.  This proves (4.2).  Equation (3.2) now gives
\[
  L_-(z)-Y(z)=\sum_{s\geq2}\frac{Y(z)^s}{s},
\]
whose coefficients are nonnegative because those of $Y$ are nonnegative.
Finally, if $m$ is odd, then $m-1$ is even, so
$(-\lambda)^{m-1}=\lambda^{m-1}$, which gives (4.3).
\end{proof}

\begin{remark}[What is, and is not, being discarded]
The coefficient $m\coeff{m}{L_--Y}$ is precisely the aggregate of all
length-$m$ words with at least two excursions.  The lemma proves that this
aggregate is nonnegative, so it may be removed from a lower bound.  It does
not say that each such projection word is nonnegative.  Similarly, the
alternating signs within a single exceptional excursion are removed only
after all possible excursion lengths are summed into $k_m(\Lambda)$.
\end{remark}

\section{The one-excursion contribution after isolating $\Lambda$}

Under the decomposition (1.2),
\[
  \mu(\{\Lambda\})=\gamma^2,
  \qquad
  \mu([-M,M])=\norm{g_s}^2,
  \tag{5.1}
\]
where
\[
  M:=\sqrt{pq-\Lambda^2}.
  \tag{5.2}
\]
The reason for the support assertion in (5.1) will be proved in
Section~6.

It is useful first to see the exceptional contribution directly from words.
Fix $1\leq b\leq m-1$.  A word with one exceptional excursion of length $b$
has two transition factors $-\gamma$, $b-1$ internal factors $-\Lambda$, and
$m-b-1$ constant factors $p$.  Therefore each cyclic rotation has weight
\[
  p^{m-b-1}\gamma^2(-\Lambda)^{b-1}.
  \tag{5.3}
\]
There are $m$ rotations, and summing over $b$ gives
\begin{equation}
\begin{aligned}
  m\gamma^2\sum_{b=1}^{m-1}p^{m-b-1}(-\Lambda)^{b-1}
  &=m\gamma^2
    \frac{p^{m-1}-(-\Lambda)^{m-1}}{p+\Lambda}\\
  &=m k_m(\Lambda)\gamma^2.
\end{aligned}  
  \tag{5.4}
\end{equation}
The last equality again uses that $m-1$ is even.

Likewise, the sum over all lengths of a single remaining-subspace excursion
is
\begin{equation}
\begin{aligned}
  &m\sum_{b=1}^{m-1}
    p^{m-b-1}\ip{g_s}{(-A_s)^{b-1}g_s}\\
  &\hspace{3em}=m\ip{g_s}{k_m(A_s)g_s},
\end{aligned}
  \tag{5.5}
\end{equation}
where $k_m(A_s)$ is defined by the polynomial identity
\[
  k_m(A_s)=\sum_{j=0}^{m-2}p^{m-2-j}(-A_s)^j.
  \tag{5.6}
\]

\begin{lemma}[Minimum of the one-excursion polynomial]
If $0\leq M<p$, then
\[
  k_m(\lambda)\geq k_m(M)
  \qquad\text{for every }\lambda\in[-M,M].
  \tag{5.7}
\]
Consequently,
\[
  m\coeff{m}{L_-(z)}
  \geq
  m k_m(\Lambda)\gamma^2
  +m k_m(M)\norm{g_s}^2.
  \tag{5.8}
\]
\end{lemma}

\begin{proof}
For $0\leq s<p$,
\[
  k_m(-s)=\frac{p^{m-1}-s^{m-1}}{p-s}
  >\frac{p^{m-1}-s^{m-1}}{p+s}=k_m(s).
  \tag{5.9}
\]
It therefore suffices to minimize on $[0,M]$.  Put $n=m-1$.  For
$0\leq\lambda<p$,
\[
  k_m'(\lambda)
  =-\frac{p^n+np\lambda^{n-1}+(n-1)\lambda^n}
          {(p+\lambda)^2}<0.
  \tag{5.10}
\]
Thus $k_m$ is decreasing on $[0,M]$, proving (5.7).  Combining (4.3),
(5.1), and (5.7) gives (5.8).
\end{proof}

This proves, in projection-word language, the origin of the two terms
\[
  m k_m(\Lambda)\gamma^2,
  \qquad
  m k_m(M)\norm{g_s}^2.
\]
They are the lower bound obtained from all one-excursion words.  No
interpretation is assigned to the two scalar products beyond this exact
enumeration.

\section{The two words with no return to $P_0$}

It remains to control the words $\D^m$ and $\Ssafe^m$.  The required estimate
comes from the following Hilbert--Schmidt bound.

\begin{lemma}[Hilbert--Schmidt bound]
One has
\[
  \Tr(A^2)+2\norm{g}^2\leq pq.
  \tag{6.1}
\]
Consequently, there is at most one eigenvalue of $A$ larger than $q$.  If
$\Lambda>q$ is that eigenvalue and $M$ is given by (5.2), then every other
eigenvalue $\lambda$ of $A$ satisfies
\[
  |\lambda|\leq M<q<\Lambda<\frac12<p.
  \tag{6.2}
\]
\end{lemma}

\begin{proof}
Since $0\leq U\leq1$,
\[
  \Tr(T_U^2)=\int_{[0,1]^2}U(x,y)^2\,dx\,dy
  \leq\int_{[0,1]^2}U(x,y)\,dx\,dy=q.
\]
Expanding the square of the block matrix in (1.1) gives
\[
  \Tr(T_U^2)=q^2+2\norm{g}^2+\Tr(A^2).
\]
Subtracting $q^2$ proves (6.1).

If two eigenvalues of $A$ were larger than $q$, then
\[
  \Tr(A^2)>2q^2>pq,
\]
where $2q>p$ is equivalent to $q>1/3$.  This contradicts (6.1).  After
isolating $\Lambda$, (6.1) gives
\[
  \sum_{\lambda_i\neq\Lambda}\lambda_i^2
  \leq pq-\Lambda^2-2\norm{g}^2
  \leq M^2,
\]
so every remaining eigenvalue has absolute value at most $M$.  Moreover,
\[
  M^2<pq-q^2=q(1-2q)<q^2
\]
because $\Lambda>q$ and $q>1/3$.  Hence $M<q$.  Finally,
$\Lambda^2\leq\Tr(A^2)\leq pq<1/4$, so $\Lambda<1/2<p$.
\end{proof}

The no-return contribution is
\[
  -\Tr(A^m)=-\Lambda^m-\Tr(A_s^m).
  \tag{6.3}
\]
For every eigenvalue $\lambda\in[-M,M]$ of $A_s$,
\[
  \lambda^m\leq M^{m-2}\lambda^2.
  \tag{6.4}
\]
For nonnegative $\lambda$ this is $\lambda^{m-2}\leq M^{m-2}$; for negative
$\lambda$ the left side is negative and the right side is nonnegative.
Therefore
\begin{equation}
\begin{aligned}
  -\Tr(A^m)
  &\geq-\Lambda^m-M^{m-2}\Tr(A_s^2)\\
  &\geq-\Lambda^m-M^{m-2}
    \left(M^2-2\norm{g}^2\right)\\
  &=-\Lambda^m-M^m
    +2M^{m-2}\left(\gamma^2+\norm{g_s}^2\right).
\end{aligned}
  \tag{6.5}
\end{equation}
Thus the word $\D^m$ remains as the exact term $-\Lambda^m$.  The entire
word $\Ssafe^m$ is bounded using the squared eigenvalues of $A_s$; the amount
of the Hilbert--Schmidt bound occupied by $g$ produces the last term in
(6.5).

For completeness, if $A$ has no eigenvalue larger than $q$, then
\[
  \lambda^m\leq q^{m-2}\lambda^2
\]
for every eigenvalue $\lambda$ of $A$.  Hence
\[
  \Tr(A^m)
  \leq q^{m-2}\Tr(A^2)
  \leq q^{m-2}\bigl(\Tr(A^2)+2\norm{g}^2\bigr)
  \leq pq^{m-1}.
  \tag{6.6}
\]
Since $\coeff{m}{L_-}\geq0$, identity (3.4) already gives the desired cycle
bound.  This explains why Proposition~5.4 only needs to treat the single
exceptional direction.

\section{Assembly: Proposition~5.4}

Define
\[
  \Lambda_m:=2M^{m-2}+mk_m(\Lambda),
  \qquad
  M_m:=2M^{m-2}+mk_m(M),
  \tag{7.1}
\]
and
\[
  R_m:=\Lambda^m+M^m-pq^{m-1}.
  \tag{7.2}
\]

\begin{proposition}[Excursion form of Proposition~5.4]
Under the assumptions above,
\[
\boxed{
  t(C_m,W)-\bigl(p^m-pq^{m-1}\bigr)
  \geq
  -R_m+\Lambda_m\gamma^2+M_m\norm{g_s}^2.
}
  \tag{7.3}
\]
Moreover, $\Lambda_m>0$ and $M_m>0$.
\end{proposition}

\begin{proof}
Subtract the proposed lower bound from the exact trace identity (3.4):
\begin{equation}
\begin{aligned}
  t(C_m,W)-\bigl(p^m-pq^{m-1}\bigr)
  &=pq^{m-1}-\Tr(A^m)+m\coeff{m}{L_-(z)}.
\end{aligned}
  \tag{7.4}
\end{equation}
The words with no visit to $P_0$ satisfy (6.5), while the complete aggregate
of words that visit $P_0$ satisfies (5.8).  Substituting both estimates into
(7.4) gives
\[
\begin{aligned}
  &t(C_m,W)-\bigl(p^m-pq^{m-1}\bigr)\\
  &\quad\geq
  pq^{m-1}-\Lambda^m-M^m
  +2M^{m-2}\bigl(\gamma^2+\norm{g_s}^2\bigr)\\
  &\qquad\quad
  +mk_m(\Lambda)\gamma^2
  +mk_m(M)\norm{g_s}^2\\
  &\quad=
  -\bigl(\Lambda^m+M^m-pq^{m-1}\bigr)\\
  &\qquad\quad
  +\bigl(2M^{m-2}+mk_m(\Lambda)\bigr)\gamma^2
  +\bigl(2M^{m-2}+mk_m(M)\bigr)\norm{g_s}^2,
\end{aligned}
\]
which is (7.3).  Finally, $0\leq M<\Lambda<p$ and
\[
  k_m(x)=\frac{p^{m-1}-x^{m-1}}{p+x}>0
  \qquad(0\leq x<p),
\]
so both coefficients in (7.1) are positive.
\end{proof}

\section{Summary of the bookkeeping}

The proof can be read directly from the following decomposition:
\[
\begin{array}{rcl}
\text{all-$\R$ word}
  &:& p^m,\\[2mm]
\text{no-return exceptional word}
  &:& -\Lambda^m,\\[2mm]
\text{no-return remaining-subspace word}
  &:& -\Tr(A_s^m),\\[2mm]
\text{all one-excursion words}
  &:& m\displaystyle\int k_m(\lambda)\,d\mu(\lambda),\\[3mm]
\text{all words with at least two excursions}
  &:& m\coeff{m}{(L_-(z)-Y(z))}\geq0.
\end{array}
\tag{8.1}
\]
The first line cancels when $p^m$ is subtracted.  The last line may be
discarded only after aggregation.  Splitting the fourth line at the atom
$\Lambda$ gives $mk_m(\Lambda)\gamma^2$ and the lower bound
$mk_m(M)\norm{g_s}^2$.  Finally, the Hilbert--Schmidt estimate controls the
third line and supplies $2M^{m-2}(\gamma^2+\norm{g_s}^2)$.  These are exactly
the terms appearing in Proposition~5.4.

\end{document}
